% =====================================================================
%  1bat-ud3-11.tex  —  HORA 11: Representar funcions a trossos
%  (sense preàmbul: s'inclou des de main.tex)
% =====================================================================
\horatitol{Sessió 11 — Representar funcions a trossos}%
{La notació amb claus de sistema, la representació amb punts oberts i tancats, i el tancament matemàtic del debat fiscal: la continuïtat.}

\subsection{Idea clau}

A la sessió anterior vau descobrir el «salt» injust del sistema fiscal ingenu. Avui
li posem nom i el resolem. Una \textbf{funció definida a trossos} és una funció que
té \textbf{expressions diferents en intervals diferents}. S'escriu amb una
\textbf{clau de sistema}, indicant per a cada tros la seva fórmula i el seu interval.

\begin{idea}
\textbf{Funció a trossos:} cada tram té expressió i interval. Exemple:
\[
f(x)=
\begin{cases}
2x+1 & \text{si } x<2,\\[2pt]
x+3 & \text{si } x\ge 2.
\end{cases}
\]
\textbf{Punts oberts i tancats a les fronteres:}
\begin{itemize}[leftmargin=1.4em, itemsep=2pt]
  \item Un \textbf{punt tancat} ($\bullet$) indica que aquell valor \textbf{sí} que
        pertany al tros (desigualtat $\le$ o $\ge$).
  \item Un \textbf{punt obert} ($\circ$) indica que aquell valor \textbf{no} pertany
        al tros (desigualtat $<$ o $>$).
\end{itemize}
\textbf{Continuïtat:} si en una frontera els dos trossos \emph{s'enganxen} (acaben i
comencen al mateix punt), la funció és \textbf{contínua} (no fa salt). Si queden a
altures diferents, hi ha un \textbf{salt} (discontinuïtat).
\end{idea}

\subsection{Exemples}

\begin{exemple}[una funció a trossos contínua]
Representem
\[
f(x)=
\begin{cases}
x+1 & \text{si } x\le 2,\\[2pt]
-x+5 & \text{si } x>2.
\end{cases}
\]
A la frontera $x=2$: el primer tros arriba a $f(2)=3$ (punt \textbf{tancat}, perquè
$x\le 2$); el segon tros, just després de $2$, també surt de $3$ (punt \textbf{obert}
en $x=2$, però amb el mateix valor). Com que tots dos coincideixen en $3$, la funció
és \textbf{contínua}.

\begin{center}
\begin{tikzpicture}
\begin{axis}[udaxis, width=9cm, height=5.2cm,
    xlabel={\small $x$}, ylabel={\small $y$},
    xmin=-2.5, xmax=6, ymin=-1, ymax=5,
    xtick={-2,0,2,4,6}, ytick={0,1,3,5}]
  \addplot[udblue, domain=-2:2, samples=2]{x+1};
  \addplot[udblue, domain=2:5, samples=2]{-x+5};
  \addplot[udnavy, only marks, mark=*, mark size=1.8pt] coordinates {(2,3)};
  \node[udnavy, font=\scriptsize, anchor=south west] at (axis cs:2,3) {$(2,3)$};
\end{axis}
\end{tikzpicture}

\figcap{Funció contínua: els dos trossos s'enganxen al punt $(2,3)$, sense salt.}
\end{center}
\end{exemple}

\begin{exemple}[una funció a trossos amb salt]
Ara
\[
g(x)=
\begin{cases}
x+1 & \text{si } x\le 2,\\[2pt]
x+3 & \text{si } x>2.
\end{cases}
\]
A $x=2$: el primer tros acaba (tancat) a $g(2)=3$; el segon tros comença (obert) just
després de $2$ a un valor proper a $5$. Els dos trossos \textbf{no s'enganxen}: hi ha
un \textbf{salt} de $3$ a $5$.

\begin{center}
\begin{tikzpicture}
\begin{axis}[udaxis, width=9cm, height=5.2cm,
    xlabel={\small $x$}, ylabel={\small $y$},
    xmin=-2.5, xmax=6, ymin=-1, ymax=9,
    xtick={-2,0,2,4,6}, ytick={0,3,5,7}]
  \addplot[udblue, domain=-2:2, samples=2]{x+1};
  \addplot[udblue, domain=2:5, samples=2]{x+3};
  \addplot[udnavy, only marks, mark=*, mark size=1.8pt] coordinates {(2,3)};
  \addplot[udnavy, only marks, mark=o, mark size=1.8pt] coordinates {(2,5)};
  \draw[udgrey, dashed] (axis cs:2,3) -- (axis cs:2,5);
  \node[udnavy, font=\scriptsize, anchor=east] at (axis cs:1.9,4) {salt};
\end{axis}
\end{tikzpicture}

\figcap{Funció amb salt: punt tancat $(2,3)$ i punt obert $(2,5)$; la gràfica es trenca.}
\end{center}
\end{exemple}

\subsection{Exercicis resolts}

\begin{exresolt}[avaluar una funció a trossos]
Donada
\[
f(x)=
\begin{cases}
-2x & \text{si } x<0,\\[2pt]
x+1 & \text{si } 0\le x<3,\\[2pt]
4 & \text{si } x\ge 3,
\end{cases}
\]
calcula $f(-2)$, $f(0)$, $f(2)$ i $f(3)$.
\solucio
Triem el tros segons on cau cada $x$:
\begin{itemize}[leftmargin=1.4em, itemsep=1pt]
  \item $f(-2)$: com que $-2<0$, fem servir $-2x$: $f(-2)=-2\cdot(-2)=4$.
  \item $f(0)$: com que $0\le 0<3$, fem servir $x+1$: $f(0)=0+1=1$.
  \item $f(2)$: com que $0\le 2<3$, fem servir $x+1$: $f(2)=2+1=3$.
  \item $f(3)$: com que $3\ge 3$, fem servir $4$: $f(3)=4$.
\end{itemize}
\resposta{$f(-2)=4$, \quad $f(0)=1$, \quad $f(2)=3$, \quad $f(3)=4$.}
\end{exresolt}

\begin{exresolt}[el tancament del cas fiscal: la quota íntegra]
Expliquem com s'arregla el sistema fiscal de la sessió 10 perquè sigui \emph{just}
(continu). Amb els trams $10\%$ (fins a $\num{15000}$) i $20\%$ (de $\num{15000}$ a
$\num{30000}$), l'impost \textbf{per trams} (quota íntegra) és:
\[
I(x)=
\begin{cases}
0,10\,x & \text{si } 0\le x\le \num{15000},\\[2pt]
\num{1500}+0,20\,(x-\num{15000}) & \text{si } \num{15000}< x\le \num{30000}.
\end{cases}
\]
\solucio
La clau és que en el segon tram el $20\%$ s'aplica \textbf{només a la part que passa
de $\num{15000}$}, i s'hi suma el que ja s'havia pagat pel primer tram
($0,10\cdot\num{15000}=\num{1500}\eur$). Comprovem la frontera $x=\num{15000}$:
\begin{itemize}[leftmargin=1.4em, itemsep=1pt]
  \item Primer tram: $I(\num{15000})=0,10\cdot\num{15000}=\num{1500}\eur$.
  \item Segon tram, just a $\num{15000}$: $\num{1500}+0,20\cdot 0=\num{1500}\eur$.
\end{itemize}
Coincideixen: la funció \textbf{s'enganxa} a $\num{15000}$, és \textbf{contínua} i ja
\textbf{no hi ha salt}. Guanyar una mica més de renda ja no fa baixar el net.
\resposta{Aplicant el \% només a la part de cada tram (quota íntegra), els trossos
s'enganxen a la frontera: la funció és contínua i justa.}
\end{exresolt}

\subsection{Exercicis per fer}

\begin{exercicis}
\begin{exlist}
  \item Per a la funció
        \[
        f(x)=
        \begin{cases}
        x+2 & \text{si } x<1,\\[2pt]
        3x & \text{si } x\ge 1,
        \end{cases}
        \]
        calcula $f(-3)$, $f(0)$, $f(1)$ i $f(4)$.
  \item Representa en paper mil·limetrat la funció
        \[
        f(x)=
        \begin{cases}
        2 & \text{si } x<0,\\[2pt]
        x+2 & \text{si } x\ge 0,
        \end{cases}
        \]
        i marca amb cura el punt obert i el punt tancat a la frontera $x=0$. És
        contínua?
  \item Mira la gràfica i digues si la funció és contínua o té un salt a la frontera.
        Indica quin punt és obert i quin és tancat.

        \begin{center}
        \begin{tikzpicture}
        \begin{axis}[udaxis, width=8.6cm, height=4.8cm,
            xlabel={\small $x$}, ylabel={\small $y$},
            xmin=-3, xmax=5, ymin=-1, ymax=6,
            xtick={-2,0,2,4}, ytick={0,2,4}]
          \addplot[udblue, domain=-2:1, samples=2]{0.5*x+1};
          \addplot[udblue, domain=1:4, samples=2]{x+2};
          \addplot[udnavy, only marks, mark=*, mark size=1.7pt] coordinates {(1,1.5)};
          \addplot[udnavy, only marks, mark=o, mark size=1.7pt] coordinates {(1,3)};
          \draw[udgrey, dashed] (axis cs:1,1.5) -- (axis cs:1,3);
        \end{axis}
        \end{tikzpicture}

        \figcap{Llegeix la frontera: punt obert, punt tancat i continuïtat.}
        \end{center}

  \item Un ajut social es concedeix per trams segons els ingressos mensuals $x$:
        \[
        A(x)=
        \begin{cases}
        300 & \text{si } x<500,\\[2pt]
        200 & \text{si } 500\le x<1000,\\[2pt]
        0 & \text{si } x\ge 1000.
        \end{cases}
        \]
        Calcula l'ajut per a $x=400$, $x=500$, $x=900$ i $x=1000$, i digues si la
        funció té salts (i a on).
  \item \textbf{(Continuïtat)} Per a la funció
        \[
        f(x)=
        \begin{cases}
        x+1 & \text{si } x\le 2,\\[2pt]
        2x+k & \text{si } x>2,
        \end{cases}
        \]
        troba el valor de $k$ perquè la funció sigui \textbf{contínua} a $x=2$.
        (Pista: els dos trossos han de donar el mateix valor a $x=2$.)
  \item \textbf{(Interpretació)} A l'ajut social de l'exercici 4, què signifiquen els
        salts en el context real? És just que algú que cobra $999\eur$ rebi $200\eur$
        i algú que cobra $\num{1000}\eur$ no rebi res? Relaciona-ho amb el debat fiscal
        de la sessió anterior.
\end{exlist}
\end{exercicis}

% ---------------------------------------------------------------------
\fullsolucions{Sessió 11}
\begin{solucions}
\begin{sollist}
  \item $f(-3)=-3+2=-1$ (tros $x+2$); $f(0)=0+2=2$ (tros $x+2$, perquè $0<1$);
        $f(1)=3\cdot 1=3$ (tros $3x$, perquè $1\ge 1$); $f(4)=3\cdot 4=12$.
  \item A $x=0$: el tros $2$ (per $x<0$) deixa un punt \textbf{obert} a $(0,2)$; el
        tros $x+2$ (per $x\ge 0$) dóna un punt \textbf{tancat} a $(0,2)$. Com que tots
        dos valen $2$, \textbf{sí que és contínua} (s'enganxen).
  \item Hi ha un \textbf{salt}: punt tancat $(1;\,1,5)$ (tros de l'esquerra, $x\le 1$)
        i punt obert $(1,3)$ (tros de la dreta). No és contínua.
  \item $A(400)=300$; $A(500)=200$; $A(900)=200$; $A(1000)=0$. Té salts a $x=500$
        (de $300$ a $200$) i a $x=\num{1000}$ (de $200$ a $0$).
  \item A $x=2$: tros esquerre $f(2)=2+1=3$; tros dret $2\cdot 2+k=4+k$. Perquè
        coincideixin: $4+k=3 \Rightarrow k=-1$.
  \item Els salts signifiquen que un canvi mínim d'ingressos fa perdre de cop una part
        de l'ajut (o tot). No és just que $1\eur$ més de sou faci perdre $200\eur$
        d'ajut: és el mateix problema dels «salts» fiscals: caldria un disseny continu
        (l'ajut que disminueixi gradualment) perquè ningú surti perjudicat per guanyar
        una mica més.
\end{sollist}
\end{solucions}
